% =====================================================================
%  1bat-ud2-2.tex  —  HORA 2: La regla general d'un servei (descoberta)
%  (sense preàmbul: s'inclou des de main.tex)
% =====================================================================
\horatitol{Sessió 2 — La regla general d'un servei}%
{Deduir una expressió general (cost fix + cost variable) a partir de dades reals.}

\subsection{Idea clau}

Moltes situacions de l'àmbit social tenen una part \textbf{fixa} (no canvia) i una part
\textbf{variable} (depèn de quantes persones, hores o unitats). La regla general
s'escriu sempre amb la mateixa estructura:
\[
\boxed{\ \text{cost total} \;=\; \underbrace{\text{cost fix}}_{\text{terme independent}}
\;+\; \underbrace{\text{cost per unitat}}_{\text{coeficient}} \per x\ }
\]
on $x$ és el nombre d'unitats (usuaris, hores, àpats\ldots).

\subsection{Exemples}

\begin{exemple}[centre de dia]
Un centre de dia té $\num{2000}\eur$ de despeses fixes mensuals (lloguer i
subministraments) i, a més, gasta $15\eur$ per cada usuari atès (àpats i material).
El cost total mensual per a $x$ usuaris és
\[
C = \num{2000} + 15x \quad(\text{en euros}).
\]
La part fixa $\num{2000}$ hi és sempre; la part $15x$ creix amb el nombre d'usuaris.
\end{exemple}

\begin{exemple}[llegir la regla]
A l'expressió $C=\num{2000}+15x$:
\begin{itemize}[leftmargin=1.4em, itemsep=2pt]
  \item El \textbf{terme independent} és $\num{2000}$: el cost si no hi hagués cap usuari.
  \item El \textbf{coeficient} de $x$ és $15$: el que augmenta el cost per cada usuari nou.
\end{itemize}
\end{exemple}

\subsection{Exercicis resolts}

\begin{exresolt}[construir la regla general]
Un servei de transport adaptat cobra una quota fixa de $\num{30}\eur$ al mes,
més $\num{1,5}\eur$ per cada trajecte. Escriu el cost mensual per a $x$ trajectes
i calcula'l per a $x=20$ trajectes.
\solucio
Part fixa: $30$. Part variable: $\num{1,5}$ per trajecte, és a dir $\num{1,5}\,x$.
\[
C = 30 + \num{1,5}\,x \quad(\text{en euros}).
\]
Per a $x=20$:
\[
C = 30 + \num{1,5}\per 20 = 30 + 30 = 60\eur.
\]
\resposta{$C = 30 + \num{1,5}\,x$; \; per a $20$ trajectes, $C=60\eur$.}
\end{exresolt}

\begin{exresolt}[de la taula a l'expressió]
En una ludoteca s'han anotat aquests costos segons el nombre d'infants:

\begin{center}
\begin{tabular}{c c}
\toprule
Infants ($x$) & Cost ($\eur$) \\
\midrule
0 & 500 \\
1 & 508 \\
2 & 516 \\
3 & 524 \\
\bottomrule
\end{tabular}
\end{center}

Troba la regla general que dona el cost en funció de $x$.
\solucio
Quan $x=0$, el cost és $500$: aquest és el \textbf{terme independent}.
Cada infant nou afegeix $508-500=8\eur$ (i es manté: $516-508=8$, $524-516=8$):
el \textbf{coeficient} de $x$ és $8$. Per tant:
\[
C = 500 + 8x.
\]
Comprovació amb $x=3$: $C=500+8\per 3 = 524\eur$. \checkmark
\resposta{$C = 500 + 8x$}
\end{exresolt}

\subsection{Exercicis per fer}

\begin{exercicis}
\begin{exlist}
  \item Una associació lloga un local per $\num{200}\eur$ i, a més, paga $5\eur$
        de material per cada participant d'un taller. Escriu el cost total per a
        $x$ participants.
  \item Un menjador social té $\num{1200}\eur$ de despeses fixes i gasta $\num{3,5}\eur$
        per àpat servit. Escriu el cost mensual per a $x$ àpats i calcula'l per a
        $x=400$ àpats.
  \item A partir de la taula, troba la regla general $C$ en funció de $x$:
    \begin{center}
    \begin{tabular}{c c}
    \toprule
    $x$ & $C$ ($\eur$) \\
    \midrule
    0 & 150 \\
    1 & 162 \\
    2 & 174 \\
    \bottomrule
    \end{tabular}
    \end{center}
  \item Una entitat rep una subvenció fixa de $\num{1000}\eur$ i, a més, ingressa
        $20\eur$ per cada soci. Escriu els ingressos totals si té $x$ socis.
  \item Identifica, a cada expressió, el terme independent i el coeficient de $x$:
        \quad (a) $C = 800 + 12x$ \quad (b) $C = 45x$ \quad (c) $C = 250 + \num{2,5}x$.
  \item \textbf{(Raonament)} Dos serveis tenen els costos $C_1 = 100 + 10x$ i
        $C_2 = 300 + 6x$. Sense calcular gaire, explica quin servei surt més barat
        quan hi ha \emph{pocs} usuaris i per què la part fixa hi té un paper important.
\end{exlist}
\end{exercicis}

% ---------------------------------------------------------------------
\fullsolucions{Sessió 2}
\begin{solucions}
\begin{sollist}
  \item $C = 200 + 5x$ (euros)
  \item $C = 1200 + \num{3,5}\,x$; \; per a $x=400$: $C = \num{2600}\eur$
  \item $C = 150 + 12x$
  \item Ingressos $= 1000 + 20x$ (euros)
  \item (a) terme independent $800$, coeficient $12$ \quad
        (b) terme independent $0$, coeficient $45$ \quad
        (c) terme independent $250$, coeficient $\num{2,5}$
  \item Amb pocs usuaris surt més barat $C_1$ (té la part fixa més baixa, $100$
        davant de $300$); quan $x$ és petit, el terme variable pesa poc i mana
        la part fixa.
\end{sollist}
\end{solucions}
